% !TEX program = xelatex
% 全局配置：只改本文件的软件名、版本号和字体，其余章节会自动引用。
% ============================================================

% Document info —— 改成你自己的软件信息
\newcommand{\Name}{DemoTool}
\newcommand{\NameCN}{轻量数据整理工具}
\newcommand{\FullName}{Lightweight Data Cleaning Toolkit}
\newcommand{\Version}{V1.0}
\newcommand{\Header}{\NameCN}

% 西文字体：正文 Times New Roman，无衬线 Arial，代码 DejaVu Sans Mono。
% 中文字体由 easybook.cls / ctex 管理，常用选项：font = noto 或 font = windows。
\newcommand{\ManualRomanFont}{Times New Roman}
\newcommand{\ManualSansFont}{Arial}
\newcommand{\ManualMonoFont}{DejaVu Sans Mono}
\AtBeginDocument{%
  \setmainfont{\ManualRomanFont}%
  \setsansfont{\ManualSansFont}%
  \setmonofont{\ManualMonoFont}%
}

% 代码框外观开关：手册用 decorated（默认带边框底色），源程序文档在
% source-code.tex 里改为 plain（纯白底无边框）。一般无需改这里。
\providecommand{\CodeBlockStyle}{decorated}
\newcommand{\PlainCodeBlockStyle}{plain}
\ifx\CodeBlockStyle\PlainCodeBlockStyle
  \lstdefinestyle{blockstyle}{
    style=basecode,
    basicstyle=\ttfamily\footnotesize,
    numbers=none,
    xleftmargin=1em,
    framexleftmargin=0pt,
    backgroundcolor={},
    frame=none
  }
  \lstdefinestyle{codestyle}{
    style=basecode,
    basicstyle=\ttfamily\scriptsize,
    numbers=none,
    numbersep=10pt,
    numberstyle=\tiny\color{codegray},
    xleftmargin=1em,
    framexleftmargin=0pt,
    backgroundcolor={},
    frame=none
  }
\fi

\usepackage{tikz}
\usetikzlibrary{positioning, arrows.meta, calc, fit}

\usepackage{zhlipsum,bicaption,tcolorbox}
\usepackage{longtable}
\usepackage{booktabs}
\usepackage{array}
\usepackage{ragged2e}
\usepackage{multirow}
\usepackage{float}
\usepackage[section]{placeins}
\usepackage[version=4]{mhchem}

% 定宽表格列：左对齐且禁止英文单词在单元格中断词。
\newcolumntype{L}[1]{>{\RaggedRight\arraybackslash\hyphenpenalty=10000\exhyphenpenalty=10000}p{#1}}

% 列表符号和正文一起缩进。
\setlist[itemize,1]{leftmargin=*, labelindent=2em, itemsep=2pt, topsep=4pt, parsep=0pt}
\setlist[enumerate,1]{leftmargin=*, labelindent=2em, itemsep=2pt, topsep=4pt, parsep=0pt}

% 行内长标识符在行尾按字符折行，避免把上一行中文撑开。
\newcommand{\cbr}[1]{\texttt{\seqsplit{#1}}}

% 放宽断行容忍度，避免长命令或路径把两端对齐的中文行撑开。
\tolerance=10000
\emergencystretch=3em
\hbadness=10000
